\chapter{Perspectivas y trabajo a futuro}
\label{cap:perspectivas}
%% ============================================================================
%% BORRADOR - Capítulo 6
%% Los comentarios que empiezan con "%% CAMI" son notas para revisar.
%% ============================================================================

Las limitaciones identificadas en el Capítulo \ref{cap:discusion}, en particular la baja potencia estadística y la dificultad de
separar los grupos con métricas globales, orientan el trabajo a futuro en dos líneas
complementarias: una que se encuentra actualmente en desarrollo, la modelación con redes
bayesianas, y otra de más largo plazo, la puesta en funcionamiento de un simulador de datos
metagenómicos.

\section{Redes bayesianas (en desarrollo)}

Los índices de diversidad resumen toda la comunidad en un solo número y las pruebas de hipótesis
comparan esos resúmenes entre grupos; ninguno de los dos enfoques describe cómo se relacionan los
microorganismos entre sí ni cómo esas relaciones se vinculan con el estado de la planta. Para
abordar esta pregunta se está desarrollando un modelo de redes bayesianas, un modelo gráfico
probabilístico en el que un grafo acíclico dirigido $G = (V, E)$ codifica las dependencias
condicionales entre variables, de modo que la distribución conjunta se factoriza como
\begin{equation}
P(X_1, \dots, X_n) = \prod_{i=1}^{n} P\left(X_i \mid \mathrm{Pa}(X_i)\right),
\label{eq:factorizacion-bn}
\end{equation}
donde $\mathrm{Pa}(X_i)$ denota el conjunto de padres del nodo $X_i$ en el grafo.

En este trabajo, los nodos de la red son la abundancia del género fúngico \textit{Fusarium}, el
índice de Shannon como medida de diversidad alfa, la abundancia de los taxones más abundantes y el
estado de salud de la planta. La estructura se aprende a partir de los datos, por ejemplo con el
paquete \texttt{bnlearn} de R \citep{scutari2010bnlearn}, y su robustez se evalúa mediante
\emph{bootstrapping}, conservando solo los arcos que aparecen en una proporción alta de las redes
remuestreadas, y validación cruzada de $k$ particiones para medir la capacidad predictiva del
modelo.
%% CAMI: actualizar esta sección con lo que ya tengas: algoritmo de aprendizaje de estructura
%% (hc, tabu, pc...), discretización o no de las abundancias, número de taxones, umbral de
%% fuerza de arcos y primeros resultados.

Este enfoque complementa los resultados de esta tesis en tres aspectos:
\begin{itemize}
    \item Permite pasar de preguntar \emph{si} los grupos difieren a preguntar \emph{qué}
    variables están relacionadas con el estado de la planta y a través de qué otras variables.
    \item Da continuación natural al hallazgo sobre la comunidad eucariota, al modelar
    explícitamente el papel de \textit{Fusarium} y su relación con la diversidad y con los taxones
    dominantes.
    \item Extiende el análisis de redes de coocurrencia exploratorio, basado en correlaciones y
    herramientas como MicNet \citep{favila2022micnet}, hacia un modelo probabilístico con el que
    se puede hacer inferencia y predicción.
\end{itemize}

Cabe advertir que, con 53 muestras, el aprendizaje de la estructura de una red bayesiana es
inestable cuando el número de nodos es grande; de ahí la importancia del \emph{bootstrapping} y de
limitar la red a un número reducido de taxones. Esta restricción conecta directamente con la
segunda línea de trabajo.

\section{Simulador de datos metagenómicos}

A más largo plazo, se propone contar con un simulador de datos metagenómicos plenamente funcional,
que permita generar conjuntos de datos adicionales para fortalecer el análisis y llegar a
resultados más significativos. Como parte de este proyecto se desarrolló una primera versión de
PyMetaSeem, un simulador escrito en Python que genera lecturas a partir de genomas de referencia en
formato FASTA. PyMetaSeem recibe un perfil taxonómico con la cantidad de lecturas por genoma, corta
lecturas aleatorias de longitud determinada, calcula su reverso complementario y produce archivos
FASTQ con la calidad asociada. Se concibió como una alternativa más sencilla de instalar y usar que
simuladores como CAMISIM \citep{fritz2019camisim}.

El flujo de trabajo propuesto consiste en:
\begin{enumerate}
    \item Estimar, a partir de los datos reales, los perfiles taxonómicos característicos de las
    plantas saludables y no saludables: las abundancias medias de cada taxón, su variabilidad y
    las relaciones entre taxones, incluidas las aprendidas por la red bayesiana.
    \item Generar nuevos perfiles taxonómicos muestreando de esos modelos y convertirlos, con
    PyMetaSeem, en metagenomas sintéticos, es decir, archivos de lecturas FASTQ.
    \item Procesar los metagenomas sintéticos con el mismo flujo del Capítulo 3: clasificación con
    Kraken, filtrado, índices de diversidad y pruebas de hipótesis.
\end{enumerate}

Este esquema cumple varios propósitos. En primer lugar, permite realizar análisis de potencia por
simulación: generar conjuntos de datos con distintos tamaños de muestra y tamaños de efecto
conocidos, y estimar cuántas muestras reales serían necesarias para detectar con confianza las
diferencias observadas, lo que puede orientar el diseño de futuros muestreos. En segundo lugar,
como en los datos simulados la composición verdadera es conocida, permite evaluar la exactitud del
propio flujo de análisis: qué tanto sesgan la clasificación taxonómica, la profundidad de
secuenciación desigual y la rarefacción los índices de diversidad. En tercer lugar, permite
validar las redes bayesianas: si se simulan datos a partir de una red con estructura conocida, es
posible medir qué tan bien los algoritmos de aprendizaje recuperan esa estructura con el tamaño de
muestra disponible.

Es importante precisar el alcance de esta línea de trabajo. Los datos simulados heredan los
supuestos del modelo con el que se generan, por lo que no sustituyen a nuevas muestras biológicas
como evidencia de una diferencia real entre plantas sanas y enfermas. Su valor está en cuantificar
la incertidumbre, validar los métodos y diseñar mejores experimentos. La combinación de un
simulador funcional, un modelo de redes bayesianas y nuevos muestreos dirigidos permitiría, en
conjunto, llegar a conclusiones estadísticamente más sólidas sobre la relación entre el microbioma
rizosférico y la salud de la fresa.

\section{Otras líneas}

Por último, se identifican otras extensiones naturales de este trabajo:
\begin{itemize}
    \item Aplicar pruebas de abundancia diferencial que consideren la naturaleza composicional de
    los datos, con corrección por comparaciones múltiples, para confirmar los taxones candidatos.
    \item Construir un clasificador de muestras saludables y no saludables con aprendizaje
    automático, por ejemplo un bosque aleatorio (\emph{random forest}), que además permita ordenar
    los taxones por su importancia para la clasificación.
    \item Integrar los conjuntos de datos adicionales obtenidos, con cinco categorías según el tipo
    de cultivo y el estado de la planta, para evaluar si los patrones observados se mantienen en
    otros contextos agronómicos.
\end{itemize}
